\documentclass[11pt]{article}
\usepackage[margin=1in]{geometry}
\usepackage{amsmath,amssymb,mathtools}
\usepackage{fontspec,unicode-math}
\setmainfont{Latin Modern Roman}
\setsansfont{Latin Modern Sans}
\setmathfont{Latin Modern Math}
\usepackage{xurl,hyperref}
\hypersetup{colorlinks=true,urlcolor=blue,linkcolor=blue}
\newcommand{\sref}[2]{\hyperlink{source:#1:#2}{[S#2]}}
\newcommand{\cataloguescope}[1]{\paragraph{Recorded mathematical question.}#1}
\begin{document}
% UnsolvedMath ID 9700012; source catalogue code AMR-096-0012
\setcounter{section}{9700011}
\section{Mixing of branch rotation and triangulation chains}\label{9700012}
{\small\sffamily UnsolvedMath ID 9700012 · Probability}
\subsection{Definitions and mathematical statement}

\paragraph{Shared notation.}$\mathbb N=\{1,2,\ldots\}$, and $\mathbb Z,\mathbb Q,\mathbb R,\mathbb C$ denote the integers, rationals, reals and complex numbers. Any other conventions are specified locally.


Consider the following two discrete-time finite reversible Markov chains. In the triangulation chain, a state is a triangulation of a convex regular $n$-gon. At each step choose one of its $n-3$ interior diagonals uniformly, remove it to expose a quadrilateral, and insert the other diagonal of that quadrilateral.
In the cladogram chain, states are rooted binary trees with $n$ leaves labeled $1,\ldots,n$, unlabeled internal vertices and no ordering of left and right subtrees; a stem edge above the root is included, giving $2n-1$ edges. Choose uniformly an edge other than that stem edge. Detach its descendant subtree at its upper endpoint and suppress the resulting degree-two vertex, producing an edge $e$. For each of the at most four edges adjacent to $e$, reattach the detached subtree at that edge's midpoint with probability $1/4$; any remaining probability makes no change. Both chains have uniform stationary law. If $1=\lambda_1>\lambda_2$ are their largest two transition eigenvalues, define relaxation time by $t_{\rm rel}=(1-\lambda_2)^{-1}$.
\paragraph{Statement recorded in the supplied source.}
For each of these two chains, prove that $t_{\rm rel}\le Cn^{3/2}$ for some constant $C$ independent of $n$, matching the order-$n^{3/2}$ lower bound. \sref{9700012}{2}
\subsection{Short English statement}
Prove sharp relaxation-time bounds for polygon diagonal flips and local branch rotations of labeled cladograms.
\subsection{Sources}
\begin{description}
\item[\hypertarget{source:9700012:1}{S1}] ULAM.AI and the UnsolvedMath Contributors, \emph{UnsolvedMath: A Curated Collection of Open Mathematics Problems}, record 9700012 (AMR-096-0012). CC BY 4.0. \url{https://huggingface.co/datasets/ulamai/UnsolvedMath}.
\item[\hypertarget{source:9700012:2}{S2}] David J. Aldous, \emph{Mixing Times for the Branch-Rotation Chain on Cladograms (or the Triangulation Walk)}, State-space and transition definitions, pp. 1--2, and PROBLEM, p. 2. \url{https://www.stat.berkeley.edu/~aldous/Research/OP/clad-mix.pdf}.
\end{description}
\label{end:9700012}
\end{document}
